% !TEX encoding = UTF-8 Unicode
% !TEX program = pdflatex
\documentclass[12pt,a4paper]{article}

\usepackage[utf8]{inputenc}
\usepackage[T1]{fontenc}
\usepackage[polish]{babel}
\usepackage{geometry}
\usepackage{booktabs}
\usepackage{longtable}
\usepackage{listings}
\usepackage{xcolor}
\usepackage{graphicx}
\graphicspath{{./}{docs/}}
\usepackage{parskip}
\usepackage{titlesec}
\usepackage{array}
\usepackage{caption}
\usepackage{enumitem}
\usepackage{fancyhdr}
\usepackage{lmodern}
\usepackage{amssymb}
\usepackage{float}
\usepackage{tikz}
\usetikzlibrary{shapes.geometric, positioning, arrows.meta, fit, backgrounds}
% hyperref ladowany na koncu preambury, zeby uniknac konfliktow z innymi pakietami
\usepackage{hyperref}

\geometry{
  top=2.5cm,
  bottom=2.5cm,
  left=2.5cm,
  right=2.5cm
}

% --- Kolory i styl listings ---
\definecolor{codebg}{RGB}{248,248,248}
\definecolor{codeframe}{RGB}{200,200,200}
\definecolor{keywordcolor}{RGB}{0,0,180}
\definecolor{commentcolor}{RGB}{100,100,100}
\definecolor{stringcolor}{RGB}{180,0,0}
\definecolor{accent}{RGB}{74,93,58}      % zielen marki U Alchemika

\lstdefinestyle{pythonstyle}{
  language=Python,
  backgroundcolor=\color{codebg},
  frame=single,
  rulecolor=\color{codeframe},
  basicstyle=\ttfamily\small,
  keywordstyle=\color{keywordcolor}\bfseries,
  commentstyle=\color{commentcolor}\itshape,
  stringstyle=\color{stringcolor},
  breaklines=true,
  showstringspaces=false,
  tabsize=4,
  numbers=left,
  numberstyle=\tiny\color{gray},
  captionpos=b,
  literate=%
    {ą}{{\k{a}}}1 {ć}{{\'c}}1 {ę}{{\k{e}}}1 {ł}{{\l{}}}1 {ń}{{\'n}}1
    {ó}{{\'o}}1 {ś}{{\'s}}1 {ż}{{\.z}}1 {ź}{{\'z}}1
    {Ą}{{\k{A}}}1 {Ć}{{\'C}}1 {Ę}{{\k{E}}}1 {Ł}{{\L{}}}1 {Ń}{{\'N}}1
    {Ó}{{\'O}}1 {Ś}{{\'S}}1 {Ż}{{\.Z}}1 {Ź}{{\'Z}}1
}

\lstdefinestyle{djangostyle}{
  backgroundcolor=\color{codebg},
  frame=single,
  rulecolor=\color{codeframe},
  basicstyle=\ttfamily\small,
  keywordstyle=\color{keywordcolor}\bfseries,
  commentstyle=\color{commentcolor}\itshape,
  stringstyle=\color{stringcolor},
  breaklines=true,
  showstringspaces=false,
  tabsize=2,
  captionpos=b,
  morekeywords={trans,blocktrans,load,block,extends,include,url,csrf_token,if,endif,for,endfor},
  literate=%
    {ą}{{\k{a}}}1 {ć}{{\'c}}1 {ę}{{\k{e}}}1 {ł}{{\l{}}}1 {ń}{{\'n}}1
    {ó}{{\'o}}1 {ś}{{\'s}}1 {ż}{{\.z}}1 {ź}{{\'z}}1
}

\lstdefinestyle{yamlstyle}{
  backgroundcolor=\color{codebg},
  frame=single,
  rulecolor=\color{codeframe},
  basicstyle=\ttfamily\small,
  keywordstyle=\color{keywordcolor},
  commentstyle=\color{commentcolor}\itshape,
  stringstyle=\color{stringcolor},
  breaklines=true,
  showstringspaces=false,
  tabsize=2,
  numbers=left,
  numberstyle=\tiny\color{gray},
  captionpos=b,
}

\lstdefinestyle{consolestyle}{
  backgroundcolor=\color{codebg},
  frame=single,
  rulecolor=\color{codeframe},
  basicstyle=\ttfamily\small,
  breaklines=true,
  captionpos=b,
  literate=%
    {ą}{{\k{a}}}1 {ć}{{\'c}}1 {ę}{{\k{e}}}1 {ł}{{\l{}}}1 {ń}{{\'n}}1
    {ó}{{\'o}}1 {ś}{{\'s}}1 {ż}{{\.z}}1 {ź}{{\'z}}1
}

% --- Nagłówek/stopka ---
\pagestyle{fancy}
\fancyhf{}
\rhead{U Alchemika -- Dokumentacja projektowa}
\lhead{}
\cfoot{\thepage}
\renewcommand{\headrulewidth}{0.4pt}

% --- Hyperref ---
\hypersetup{
  colorlinks=true,
  linkcolor=accent,
  urlcolor=accent,
  citecolor=accent,
  pdftitle={U Alchemika System -- Dokumentacja projektowa (Projektowanie i programowanie systemow internetowych I)},
  pdfauthor={Maciej Bulhak, Malgorzata Bulhak, Hubert Olszak},
}

% ---------------------------------------------------------------
\begin{document}

% ===================== STRONA TYTUŁOWA =====================
\begin{titlepage}
  \centering
  \vspace*{1.5cm}
  {\LARGE\bfseries Collegium Witelona \par}
  \vspace{0.5cm}
  {\large Projektowanie i programowanie systemów internetowych I \par}
  \vspace{2.5cm}

  \rule{\linewidth}{0.5mm}
  \vspace{0.6cm}
  {\Huge\bfseries U Alchemika System \par}
  \vspace{0.4cm}
  {\Large Aplikacja internetowa obiektu agroturystycznego \par}
  \vspace{0.4cm}
  {\large\itshape Dokumentacja projektowa \par}
  \rule{\linewidth}{0.5mm}

  \vspace{2cm}
  \begin{tabular}{rl}
    \textbf{Autorzy:}        & Maciej Bułhak \\[4pt]
                             & Małgorzata Bułhak \\[4pt]
                             & Hubert Olszak \\[4pt]
    \textbf{Przedmiot:}      & Projektowanie i programowanie \\
                             & systemów internetowych I \\[4pt]
    \textbf{Rok akademicki:} & 2025/2026 \\[4pt]
    \textbf{Data:}           & \today \\[4pt]
    \textbf{Stos technologiczny:} & Django 4.2 (MVT) + PostgreSQL 16 \\
                             & + Bootstrap 5 + HTMX + Docker \\
  \end{tabular}

  \vfill
  {\small Repozytorium: \url{https://github.com/Bulileusz/u-alchemika-system} \par}
\end{titlepage}

% ===================== SPIS TREŚCI =====================
\tableofcontents
\newpage

% ===================== WSTĘP =====================
\section{Wprowadzenie}

\subsection{Cel dokumentu}

Niniejszy dokument jest \textbf{dokumentacją projektową} aplikacji internetowej
\textbf{„U Alchemika”} -- witryny obiektu agroturystycznego położonego w~Górach
Izerskich. Opisuje on system od strony funkcjonalnej i~technicznej: co aplikacja
robi, jak jest zbudowana, z~jakich technologii korzysta oraz jak uruchomić ją
lokalnie i~wdrożyć na serwerze. Dokument adresowany jest do osoby, która ma
zrozumieć działanie systemu, rozwijać go lub uruchomić od zera.

\subsection{Kontekst projektu}

Aplikacja powstała w~ramach kursu \emph{Projektowanie i~programowanie systemów
internetowych~I} i~współdzieli repozytorium oraz kod z~kursem \emph{Projektowanie
systemów baz danych}. Oba przedmioty opisują ten sam system z~dwóch perspektyw.
Warstwa danych -- model relacyjny, klucze obce, indeksy, historia migracji --
została szczegółowo udokumentowana w~osobnym opracowaniu
(\texttt{docs/dokumentacja.tex}). \textbf{Niniejszy dokument koncentruje się na
warstwie internetowej}: funkcjach widocznych dla użytkownika, architekturze
aplikacji webowej oraz technologiach frontendu i~backendu. Model danych omówiono
tu skrótowo (sekcja~\ref{sec:dane}), bez powielania dokumentacji bazodanowej.

\subsection{Zespół projektowy}

\begin{longtable}{p{0.30\linewidth}p{0.62\linewidth}}
\toprule
\textbf{Osoba} & \textbf{Zakres odpowiedzialności} \\
\midrule
\endhead
Maciej Bułhak & Model danych, migracje i~integralność, warstwa dostępu do danych
  (ORM), konfiguracja projektu, konteneryzacja, dokumentacja. \\[4pt]
Małgorzata Bułhak & Moduł zapytań rezerwacyjnych (formularz, walidacja, poczta),
  panel obsługi, logowanie zdarzeń, dane obiektu i~kontakt. \\[4pt]
Hubert Olszak & Warstwa prezentacji: szablony HTML, style CSS, responsywność,
  interaktywność po stronie klienta (JavaScript), blog i~atrakcje, dane testowe. \\
\bottomrule
\end{longtable}

% ===================== CHARAKTERYSTYKA SYSTEMU =====================
\section{Charakterystyka systemu}
\label{sec:charakterystyka}

\subsection{Cel i zakres}

\textbf{U Alchemika} to witryna prezentująca ofertę noclegową obiektu i~zbierająca
\textbf{zapytania rezerwacyjne} od gości. Przyjęto świadome założenie projektowe, że
system \textbf{nie realizuje pełnej rezerwacji online} -- gość wysyła zapytanie,
a~finalizacja odbywa się poza systemem (kontakt bezpośredni lub przekierowanie do
serwisu Booking.com). Dzięki temu zakres jest spójny i~kompletny w~obrębie
zdefiniowanych przypadków użycia, zamiast obejmować pobieżnie zaimplementowany moduł
płatności i~kalendarza dostępności.

\subsection{Aktorzy}

\begin{description}[nosep]
  \item[Gość (użytkownik anonimowy)] -- przegląda ofertę (pokoje, galeria,
    atrakcje, blog), zapoznaje się z~polityką obiektu i~wysyła zapytanie
    rezerwacyjne. Nie wymaga konta ani logowania.
  \item[Obsługa / administrator] -- po uwierzytelnieniu przegląda nadesłane
    zapytania w~panelu wewnętrznym, a~całością danych (pokoje, treści, zapytania)
    zarządza przez panel Django Admin.
\end{description}

\subsection{Mapa funkcji}

Tabela~\ref{tab:strony} zestawia wszystkie trasy aplikacji wraz z~ich funkcją
i~wymaganym poziomem dostępu.

\begin{longtable}{p{0.22\linewidth}p{0.27\linewidth}p{0.34\linewidth}c}
\caption{Mapa stron aplikacji}\label{tab:strony}\\
\toprule
\textbf{Strona} & \textbf{Adres (URL)} & \textbf{Funkcja} & \textbf{Dostęp} \\
\midrule
\endfirsthead
\multicolumn{4}{c}{\tablename\ \thetable\ -- ciąg dalszy}\\
\toprule
\textbf{Strona} & \textbf{Adres (URL)} & \textbf{Funkcja} & \textbf{Dostęp} \\
\midrule
\endhead
Strona główna    & \texttt{/}                 & Hero, polecane pokoje, zajawka atrakcji, CTA kontaktu. & publiczny \\
Lista pokoi      & \texttt{/pokoje/}          & Karty aktywnych pokoi z~ceną i~pojemnością. & publiczny \\
Szczegóły pokoju & \texttt{/pokoje/<slug>/}   & Galeria, opis, udogodnienia, cena, przejście do zapytania. & publiczny \\
Galeria          & \texttt{/galeria/}         & Zdjęcia obiektu z~powiększaniem (lightbox). & publiczny \\
Atrakcje         & \texttt{/atrakcje/}        & Atrakcje turystyczne okolicy. & publiczny \\
Blog             & \texttt{/blog/}            & Lista opublikowanych wpisów. & publiczny \\
Wpis blogowy     & \texttt{/blog/<slug>/}     & Treść pojedynczego wpisu. & publiczny \\
Kontakt          & \texttt{/kontakt/}         & Dane kontaktowe obiektu (z~cache). & publiczny \\
Obiekt           & \texttt{/obiekt/}          & Polityki obiektu (doba, zwierzęta, płatności). & publiczny \\
Zapytanie        & \texttt{/zapytanie/}       & Formularz zapytania rezerwacyjnego (HTMX). & publiczny \\
Logowanie        & \texttt{/login/}           & Uwierzytelnienie obsługi. & publiczny \\
Panel zapytań    & \texttt{/panel/zapytania/} & Lista nadesłanych zapytań. & chroniony \\
Panel admina     & \texttt{/admin/}           & Pełny CRUD na wszystkich encjach. & chroniony \\
Zmiana języka    & \texttt{/i18n/setlang/}    & Przełączenie interfejsu PL$\leftrightarrow$DE. & publiczny \\
\bottomrule
\end{longtable}

\subsection{Kluczowy przepływ: wysłanie zapytania}
\label{sec:przeplyw}

Najważniejszy scenariusz biznesowy łączy kilka mechanizmów technicznych w~jednej
interakcji użytkownika:

\begin{enumerate}[nosep]
  \item Gość wypełnia formularz na stronie \texttt{/zapytanie/} (opcjonalnie
    wywołany z~karty konkretnego pokoju, który zostaje podstawiony jako wartość
    początkowa pola).
  \item Formularz jest wysyłany \textbf{asynchronicznie} (HTMX) -- bez przeładowania
    strony.
  \item Serwer waliduje dane (poprawność e-maila, kolejność i~długość dat pobytu).
  \item Po powodzeniu zapytanie zostaje \textbf{zapisane w~bazie}, na adres gościa
    trafia \textbf{e-mail potwierdzający}, a~zdarzenie zostaje \textbf{zalogowane}.
  \item Gość widzi sekcję podziękowania -- fragment HTML podmieniony w~miejscu
    formularza przez HTMX.
  \item Zapytanie pojawia się na liście w~panelu obsługi \texttt{/panel/zapytania/}.
\end{enumerate}

Ten pojedynczy przepływ angażuje jednocześnie: formularz z~walidacją, komunikację
asynchroniczną, zapis do bazy przez ORM, wysyłkę poczty oraz rejestrację zdarzeń
w~logach -- dlatego stanowi dobry punkt odniesienia dla pozostałych rozdziałów.

% ===================== ARCHITEKTURA =====================
\section{Architektura aplikacji}
\label{sec:architektura}

\subsection{Wzorzec MVT}

Aplikację zbudowano na frameworku \textbf{Django 4.2}, który realizuje wzorzec
\textbf{MVT} (Model--View--Template) -- odmianę MVC, w~której rolę kontrolera pełni
\emph{View}, a~prezentację opisuje \emph{Template}. Przepływ pojedynczego żądania
HTTP przedstawia rysunek~\ref{fig:mvt}.

\begin{figure}[H]
  \centering
  \begin{tikzpicture}[
    box/.style={rectangle, draw, thick, rounded corners, align=center,
                minimum width=2.7cm, minimum height=1.1cm, font=\small},
    lbl/.style={font=\scriptsize\itshape, midway, above},
    >={Stealth[length=2.5mm]}
  ]
    \node[box, fill=accent!10] (url)   at (0,0)    {URLconf\\\texttt{urls.py}};
    \node[box, fill=accent!10] (view)  at (3.8,0)  {View\\\texttt{views/}};
    \node[box, fill=accent!10] (model) at (7.6,1.2){Model + ORM\\\texttt{models.py}};
    \node[box, fill=accent!10] (tmpl)  at (7.6,-1.2){Template\\\texttt{templates/}};
    \node[box, fill=gray!8]    (db)    at (11.2,1.2){PostgreSQL};

    \draw[->] (-2.4,0) -- node[lbl]{żądanie} (url.west);
    \draw[->] (url) -- node[lbl]{dispatch} (view);
    \draw[->] (view) to[bend left=18] node[lbl, sloped]{zapytanie} (model);
    \draw[->] (model) to[bend left=18] node[lbl, sloped, below]{dane} (view);
    \draw[->] (model) -- node[lbl]{SQL} (db);
    \draw[->] (view) to[bend right=18] node[lbl, sloped]{kontekst} (tmpl);
    \draw[->] (tmpl) to[bend right=18] node[lbl, sloped, below]{HTML} (view);
    \draw[->] (view.south) to[bend right=22] node[lbl, below]{odpowiedź} (-2.4,-1.6);
  \end{tikzpicture}
  \caption{Przepływ żądania HTTP w~architekturze MVT}
  \label{fig:mvt}
\end{figure}

\begin{itemize}[nosep]
  \item \textbf{Model} -- definicje encji i~reguł biznesowych mapowane przez ORM
    na tabele (\texttt{apps/core/models.py}, \texttt{apps/content/models.py}).
  \item \textbf{View} -- logika obsługi żądania: pobranie danych, walidacja, wybór
    szablonu (\texttt{apps/core/views/}, \texttt{apps/content/views.py}).
  \item \textbf{Template} -- szablony HTML z~dziedziczeniem i~partialami
    (\texttt{backend/templates/} oraz katalogi \texttt{templates/} aplikacji).
\end{itemize}

\subsection{Podział na aplikacje}

Kod podzielono na dwie aplikacje Django o~rozłącznej odpowiedzialności, co
ograniczyło konflikty przy pracy zespołowej (każda osoba pracowała głównie we
własnym obszarze):

\begin{itemize}[nosep]
  \item \texttt{apps/core} -- rdzeń systemu: pokoje, udogodnienia, galeria,
    zapytania rezerwacyjne, dane obiektu i~kontakt, logi audytu,
  \item \texttt{apps/content} -- treści marketingowe: wpisy blogowe i~atrakcje
    okolicy (aplikacja niezależna, bez powiązań z~\texttt{core}).
\end{itemize}

\subsection{Struktura repozytorium}

\begin{lstlisting}[style=consolestyle, caption={Najważniejsze elementy drzewa katalogów}]
u-alchemika-system/
|- backend/
|  |- config/            # settings.py, urls.py, wsgi.py
|  |- apps/
|  |  |- core/           # pokoje, zapytania, kontakt, galeria
|  |  |  |- views/       # pakiet: rooms.py, inquiries.py
|  |  |  |- models.py forms.py urls.py admin.py
|  |  |  `- static/ templates/   # zasoby i szablony aplikacji
|  |  `- content/        # blog + atrakcje (osobna apka)
|  |- templates/         # base.html, index.html (wspolne)
|  |- static/            # style.css, obrazy
|  |- locale/de/         # katalog tlumaczen niemieckich
|  |- requirements.txt   # zaleznosci Pythona
|  `- Dockerfile
|- docs/                 # dokumentacja (ten plik + dok. bazy danych)
`- docker-compose.yml    # web + db + mailhog
\end{lstlisting}

Routing jest hierarchiczny: główny \texttt{config/urls.py} deleguje przez
\texttt{include()} do plików \texttt{urls.py} obu aplikacji oraz montuje widoki
logowania i~przełącznika języka.

% ===================== MODEL DANYCH =====================
\section{Model danych (skrót)}
\label{sec:dane}

Dane przechowuje \textbf{PostgreSQL 16}; dostęp realizuje wyłącznie \textbf{Django
ORM} -- w~kodzie nie ma ręcznie pisanego SQL. Tabela~\ref{tab:encje} wymienia
encje systemu. Pełny model relacyjny (klucze obce, reguły kasowania, indeksy,
historia migracji) opisano w~dokumentacji bazodanowej.

\begin{longtable}{p{0.24\linewidth}p{0.70\linewidth}}
\caption{Encje systemu}\label{tab:encje}\\
\toprule
\textbf{Encja} & \textbf{Rola} \\
\midrule
\endhead
\texttt{Room}        & Pokój: nazwa, slug, opis, pojemność, cena, status aktywności. \\
\texttt{RoomImage}   & Zdjęcie pokoju (relacja 1:N, kolejność wyświetlania). \\
\texttt{Amenity}     & Udogodnienie (np. WiFi, kominek). \\
\texttt{RoomAmenity} & Tabela pośrednia relacji M:N pokój--udogodnienie. \\
\texttt{Inquiry}     & Zapytanie rezerwacyjne gościa (FK do pokoju, status, daty). \\
\texttt{PropertyInfo}& Dane i~polityki obiektu (kontakt, doba hotelowa, płatności). \\
\texttt{AuditLog}    & Rejestr zdarzeń (FK do użytkownika). \\
\texttt{Post}        & Wpis blogowy (publikowany / szkic). \\
\texttt{Attraction}  & Atrakcja turystyczna okolicy. \\
\bottomrule
\end{longtable}

Optymalizację zapytań pod kątem problemu N+1 omówiono w~sekcji~\ref{sec:orm}.

% ===================== WARSTWA PREZENTACJI =====================
\section{Warstwa prezentacji}
\label{sec:frontend}

\subsection{Szablony i dziedziczenie}

Widok HTML budują szablony Django. Wspólny \texttt{base.html} zawiera nawigację,
stopkę, dołączenie Bootstrapa i~HTMX oraz blok na komunikaty; podstrony rozszerzają
go (\texttt{\{\% extends \%\}}) i~nadpisują blok treści. Powtarzalne fragmenty
wydzielono do \textbf{partiali} dołączanych przez \texttt{\{\% include \%\}} -- dzięki
temu strona główna powstaje ze złożenia niezależnych części, a~każdy autor dostarczał
własny partial bez kolizji:

\begin{lstlisting}[style=djangostyle, caption={Składanie strony głównej z partiali (index.html)}]
{% include 'core/_featured_rooms.html' %}      {# polecane pokoje #}
{% include 'content/_attractions_teaser.html' %} {# zajawka atrakcji #}
{% include 'core/_contact_cta.html' %}         {# wezwanie do kontaktu #}
\end{lstlisting}

\subsection{Style i responsywność (RWD)}

Warstwę CSS oparto na \textbf{Bootstrap 5.3} (siatka, karty, przyciski, alerty,
klasy walidacji formularza) dołączonym z~CDN, uzupełnionym o~własny arkusz
\texttt{static/css/style.css} nadający charakter marce (navbar z~logo, sekcja hero,
karty pokoi, galeria). Responsywność zapewniają: meta \texttt{viewport}, responsywna
siatka Bootstrap (karty pokoi układają się 3~w~rzędzie na desktopie i~1~na telefonie)
oraz własne media queries. Pasek nawigacji wraz z~przełącznikiem języka zawija się
na węższych ekranach.

\subsection{Interaktywność po stronie klienta}

Skrypt \texttt{apps/core/static/core/js/inquiry.js} wzbogaca formularz zapytania:
licznik znaków wiadomości, walidacja dat po stronie klienta (minimalna data~=~dziś,
data wyjazdu po dacie przyjazdu) oraz ponowne podpięcie tych zdarzeń po
asynchronicznej podmianie formularza przez HTMX (nasłuch zdarzenia
\texttt{htmx:afterSwap}). Galeria zdjęć (\texttt{gallery.html}) zawiera własny
lightbox -- powiększanie obrazu w~nakładce, zamykanej klawiszem \texttt{Esc}.

\subsection{Asynchroniczność (HTMX)}

Wysłanie formularza zapytania jest asynchroniczne dzięki \textbf{HTMX}. Formularz
deklaruje atrybuty \texttt{hx-post}, \texttt{hx-target} i~\texttt{hx-swap}, a~widok
rozpoznaje takie żądanie po nagłówku \texttt{HX-Request} i~zwraca \textbf{fragment
HTML} (a~nie JSON): przy powodzeniu -- sekcję podziękowania, przy błędach -- ten sam
formularz z~komunikatami przy polach.

\begin{lstlisting}[style=pythonstyle, caption={Rozgałęzienie odpowiedzi dla żądania HTMX (views/inquiries.py)}]
if form.is_valid():
    inquiry = form.save()
    _send_inquiry_confirmation(inquiry)
    if request.headers.get("HX-Request"):
        return render(request, "core/_inquiry_success.html", {"inquiry": inquiry})
    messages.success(request, _("Twoje zapytanie zostało wysłane!"))
    return redirect("inquiry-create")
else:
    if request.headers.get("HX-Request"):
        return render(request, "core/_inquiry_form.html", {"form": form})
\end{lstlisting}

Podejście to (\emph{progressive enhancement}) dokłada asynchroniczność do klasycznych
szablonów serwerowych bez wprowadzania ciężkiego frameworka frontendowego.

% ===================== LOGIKA APLIKACJI =====================
\section{Logika aplikacji}
\label{sec:backend}

\subsection{Formularze i walidacja}

\texttt{InquiryForm} dziedziczy po \texttt{forms.ModelForm} -- pola powstają
z~modelu \texttt{Inquiry} bez powielania definicji. Walidacja działa na kilku
poziomach: typy pól (np. \texttt{EmailField}), reguły formularza (metody
\texttt{clean\_*}) oraz reguła modelu. Każdy formularz chroni token CSRF.

\begin{lstlisting}[style=pythonstyle, caption={Walidacja wielopoziomowa (forms.py)}]
def clean_date_from(self):
    date_from = self.cleaned_data.get("date_from")
    if date_from and date_from < timezone.localdate():
        raise forms.ValidationError(_("Data przyjazdu nie może być w przeszłości."))
    return date_from

def clean(self):
    cleaned = super().clean()
    date_from, date_to = cleaned.get("date_from"), cleaned.get("date_to")
    if date_from and date_to:
        if date_to <= date_from:
            self.add_error("date_to", _("Data wyjazdu musi być późniejsza..."))
        elif (date_to - date_from).days > 30:
            self.add_error("date_to", _("Pobyt nie może trwać dłużej niż 30 dni."))
    return cleaned
\end{lstlisting}

\subsection{Dostęp do danych i optymalizacja ORM}
\label{sec:orm}

Widoki pobierają dane wyłącznie przez ORM, z~jawną optymalizacją liczby zapytań
(przeciwdziałanie problemowi N+1):

\begin{lstlisting}[style=pythonstyle, caption={Optymalizacja zapytań (views/rooms.py, views/inquiries.py)}]
# Lista pokoi: jedno dodatkowe zapytanie zamiast N (relacja 1:N)
Room.objects.filter(is_active=True).prefetch_related("images")

# Szczegóły pokoju: zdjęcia + udogodnienia (1:N oraz M:N)
Room.objects.prefetch_related("images", "amenities")

# Panel zapytań: JOIN z pokojem jednym zapytaniem (FK "w przód")
Inquiry.objects.select_related("room").order_by("-created_at")
\end{lstlisting}

\subsection{Cache}

Skonfigurowano backend \texttt{LocMemCache}. Buforowanie zastosowano w~dwóch
wariantach -- niskopoziomowym (dane kontaktowe, które rzadko się zmieniają) oraz
na poziomie całej strony (dekorator):

\begin{lstlisting}[style=pythonstyle, caption={Dwa wzorce cache (views/inquiries.py)}]
# 1. Cache niskopoziomowy - dane kontaktowe
prop = cache.get("property_info_contact")
if prop is None:
    prop = PropertyInfo.objects.first()
    cache.set("property_info_contact", prop, timeout=60 * 15)  # 15 min

# 2. Cache calej strony - dekorator
@cache_page(60 * 30)   # strona /obiekt/ buforowana na 30 minut
def property_info(request): ...
\end{lstlisting}

W~zależnościach figuruje \texttt{django-redis} jako gotowa alternatywa (\emph{drop-in})
do współdzielenia cache między procesami w~środowisku produkcyjnym.

\subsection{Poczta}

Po zapisaniu zapytania system wysyła na adres gościa e-mail potwierdzający. Treść
pochodzi z~szablonu, a~ewentualny błąd wysyłki jest przechwytywany i~logowany -- nie
przerywa zapisu zapytania. W~środowisku lokalnym SMTP kierowany jest do
\textbf{MailHog} (kontener przechwytuje pocztę; podgląd pod \texttt{http://localhost:8025}).

\begin{lstlisting}[style=pythonstyle, caption={Wysyłka potwierdzenia (views/inquiries.py)}]
def _send_inquiry_confirmation(inquiry):
    body = render_to_string("core/emails/inquiry_confirmation.txt", {"inquiry": inquiry})
    try:
        send_mail(subject=..., message=body, from_email="noreply@u-alchemika.pl",
                  recipient_list=[inquiry.email], fail_silently=False)
        logger.info("E-mail wysłany na %s", inquiry.email)
    except Exception as exc:
        logger.error("Błąd wysyłki e-maila dla zapytania #%d: %s", inquiry.pk, exc)
\end{lstlisting}

\subsection{Uwierzytelnianie i panel obsługi}

Wykorzystano wbudowany system Django: \texttt{LoginView}/\texttt{LogoutView}, model
\texttt{auth.User}, hashowanie haseł oraz walidatory ich siły. Panel zapytań chroni
dekorator \texttt{@login\_required} -- próba wejścia bez sesji przekierowuje na
\texttt{/login/} z~parametrem \texttt{next}.

\begin{lstlisting}[style=pythonstyle, caption={Ochrona panelu obsługi (views/inquiries.py)}]
@login_required(login_url="login")
def inquiry_list(request):
    inquiries = Inquiry.objects.select_related("room").order_by("-created_at")
    return render(request, "core/inquiry_list.html", {"inquiries": inquiries})
\end{lstlisting}

\subsection{Rejestrowanie zdarzeń (logger)}

W~\texttt{settings.LOGGING} skonfigurowano handler konsolowy z~formatterem oraz
logger \texttt{apps} na poziomie INFO. Rejestrowane są kluczowe zdarzenia: nowe
zapytanie (\texttt{info}), nieprawidłowy formularz (\texttt{warning}), nieudana
wysyłka poczty (\texttt{error}) oraz dostęp do panelu zapytań.

\begin{lstlisting}[style=pythonstyle, caption={Logowanie zdarzenia biznesowego (views/inquiries.py)}]
logger.info("Nowe zapytanie #%d od %s <%s> na pokój '%s' (%s-%s)",
            inquiry.pk, inquiry.full_name, inquiry.email,
            inquiry.room, inquiry.date_from, inquiry.date_to)
\end{lstlisting}

% ===================== INTERNACJONALIZACJA =====================
\section{Internacjonalizacja (PL/DE)}
\label{sec:i18n}

Interfejs jest dwujęzyczny -- polski (domyślny) i~niemiecki -- z~przełącznikiem
w~pasku nawigacji. Konfiguracja obejmuje \texttt{LANGUAGES=[pl, de]},
\texttt{LocaleMiddleware} oraz \texttt{LOCALE\_PATHS}. Przełączanie realizuje
wbudowany widok \texttt{set\_language}; wybór zapisywany jest w~cookie
\texttt{django\_language}, \textbf{bez prefiksów w~adresie} -- ścieżki pozostają
niezmienione.

\begin{lstlisting}[style=djangostyle, caption={Przełącznik języka w nawigacji (base.html)}]
<form action="{% url 'set_language' %}" method="post" class="lang-switch">
  {% csrf_token %}
  <input name="next" type="hidden" value="{{ request.path }}">
  <button name="language" value="pl">PL</button>
  <button name="language" value="de">DE</button>
</form>
\end{lstlisting}

Teksty interfejsu oznaczono \texttt{\{\% trans \%\}}/\texttt{\{\% blocktrans \%\}}
w~szablonach oraz \texttt{gettext}/\texttt{gettext\_lazy} w~kodzie Pythona (etykiety
formularza, komunikaty, walidacje). Tłumaczenia niemieckie przechowuje katalog
\texttt{backend/locale/de/LC\_MESSAGES/django.po}, kompilowany do binarnego
\texttt{.mo} przy starcie kontenera. Tłumaczeniem objęto \emph{interfejs} aplikacji;
treść redakcyjna z~bazy (opisy pokoi, wpisy blogowe) pozostaje w~języku polskim --
jest to świadome ograniczenie zakresu, opisane w~sekcji~\ref{sec:wnioski}.

% ===================== KONFIGURACJA I BEZPIECZEŃSTWO =====================
\section{Konfiguracja i bezpieczeństwo}
\label{sec:bezpieczenstwo}

\begin{itemize}[nosep]
  \item \textbf{Konfiguracja przez zmienne środowiskowe} -- dane wrażliwe (klucz,
    dostęp do bazy, \texttt{DEBUG}, \texttt{ALLOWED\_HOSTS}) wczytuje
    \texttt{python-decouple}; w~kodzie nie ma zaszytych sekretów.
  \item \textbf{Ochrona CSRF} -- token w~każdym formularzu, \texttt{CsrfViewMiddleware}.
  \item \textbf{Hasła} -- hashowanie oraz walidatory siły
    (\texttt{AUTH\_PASSWORD\_VALIDATORS}).
  \item \textbf{Ochrona przed XSS} -- automatyczne escapowanie zmiennych w~szablonach.
  \item \textbf{Clickjacking} -- \texttt{XFrameOptionsMiddleware} (nagłówek
    \texttt{X-Frame-Options}).
\end{itemize}

Zależności deklaruje \texttt{backend/requirements.txt} z~kontrolą wersji
(\texttt{Django>=4.2,<5.0}, \texttt{psycopg2-binary}, \texttt{python-decouple},
\texttt{django-redis}); instalacja odbywa się automatycznie podczas budowy obrazu
Docker, co gwarantuje powtarzalne środowisko u~każdego członka zespołu i~na serwerze.

% ===================== KONTENERYZACJA I URUCHOMIENIE =====================
\section{Konteneryzacja i uruchomienie}
\label{sec:uruchomienie}

\subsection{Środowisko (Docker Compose)}

Całość jest skonteneryzowana. \texttt{docker-compose.yml} definiuje trzy serwisy:

\begin{lstlisting}[style=yamlstyle, caption={Serwisy aplikacji (fragment docker-compose.yml)}]
services:
  db:        # PostgreSQL 16 + healthcheck + wolumen postgres_data
    image: postgres:16
    healthcheck:
      test: ["CMD-SHELL", "pg_isready -U postgres -d u_alchemika"]
  web:       # aplikacja Django
    build: ./backend
    ports: ["8000:8000"]
    depends_on:
      db:      { condition: service_healthy }
      mailhog: { condition: service_started }
  mailhog:   # przechwytywanie poczty (SMTP 1025, UI 8025)
    image: mailhog/mailhog:latest
    ports: ["1025:1025", "8025:8025"]
volumes:
  postgres_data:
\end{lstlisting}

Serwis \texttt{web} startuje dopiero po przejściu healthchecku bazy. Przy starcie
kontener wykonuje automatycznie migracje oraz kompilację tłumaczeń, po czym
uruchamia serwer deweloperski. Komunikacja między kontenerami odbywa się po
wewnętrznej sieci Dockera -- \texttt{web} łączy się z~bazą po nazwie \texttt{db}
(wewnętrzny DNS), a~na zewnątrz wystawione są jedynie zmapowane porty.

\subsection{Uruchomienie lokalne}

\begin{lstlisting}[style=consolestyle, caption={Pierwsze uruchomienie}]
# 1. Klonowanie repozytorium
git clone https://github.com/Bulileusz/u-alchemika-system.git
cd u-alchemika-system

# 2. Uruchomienie kontenerow (migracje i kompilacja tlumaczen - automatycznie)
docker compose up --build

# 3. (Nowe okno) Zaladowanie danych testowych
docker compose exec web python manage.py loaddata seed

# 4. Utworzenie konta administratora (panel i /admin/)
docker compose exec web python manage.py createsuperuser
\end{lstlisting}

Po uruchomieniu dostępne są: witryna publiczna (\texttt{http://localhost:8000/}),
panel administracyjny (\texttt{/admin/}) oraz skrzynka MailHog z~podglądem maili
(\texttt{http://localhost:8025/}). Zatrzymanie: \texttt{docker compose down}
(z~flagą \texttt{-v} także reset bazy).

\subsection{Wdrożenie zdalne}

Konteneryzacja sprawia, że wdrożenie na serwer sprowadza się do uruchomienia tego
samego \texttt{docker compose} na maszynie zdalnej. Względem konfiguracji
deweloperskiej należy wprowadzić zmiany produkcyjne zestawione w~tabeli~\ref{tab:prod}.

\begin{longtable}{p{0.30\linewidth}p{0.62\linewidth}}
\caption{Zmiany konieczne przy wdrożeniu produkcyjnym}\label{tab:prod}\\
\toprule
\textbf{Element} & \textbf{Zmiana produkcyjna} \\
\midrule
\endhead
\texttt{DEBUG} & ustawić na \texttt{False}. \\
\texttt{DJANGO\_SECRET\_KEY} & długi, losowy klucz podany przez zmienną środowiskową. \\
\texttt{DJANGO\_ALLOWED\_HOSTS} & domena/IP serwera. \\
Serwer WWW & zastąpić \texttt{runserver} serwerem produkcyjnym (np. Gunicorn) za reverse proxy (np. Nginx). \\
Pliki statyczne & \texttt{collectstatic} i~serwowanie przez Nginx (lub WhiteNoise). \\
Poczta & realny serwer SMTP zamiast MailHog (dane z~env). \\
Baza danych & trwały wolumen lub usługa zarządzana; silne hasło. \\
HTTPS & certyfikat TLS (np. Let's Encrypt) na reverse proxy. \\
\bottomrule
\end{longtable}

\noindent\textit{Uwaga: bieżąca wersja repozytorium jest skonfigurowana pod
środowisko deweloperskie (serwer \texttt{runserver}, MailHog, domyślne sekrety).
Powyższe kroki są konieczne do bezpiecznego wdrożenia produkcyjnego.}

% ===================== PODSUMOWANIE =====================
\section{Podsumowanie i decyzje projektowe}
\label{sec:wnioski}

\begin{enumerate}[leftmargin=*]
  \item \textbf{Framework porządkuje pracę.} Django pozwoliło skupić się na logice
    obiektu noclegowego zamiast na „hydraulice” (routing, ORM, formularze, panel
    admina, bezpieczeństwo). Mechanizmy gotowe z~pudełka (CSRF, hashowanie haseł,
    escapowanie XSS) podniosły bezpieczeństwo bez dodatkowego nakładu.

  \item \textbf{Konteneryzacja eliminuje problem „u~mnie działa”.} Docker Compose
    zapewnił identyczne, powtarzalne środowisko dla całego zespołu i~uprościł
    ścieżkę do wdrożenia -- ten sam plik kompozycji działa lokalnie i~na serwerze.

  \item \textbf{Technologia proporcjonalna do skali.} Zamiast ciężkiego frameworka
    frontendowego wybrano HTMX, który dokłada asynchroniczność do szablonów
    serwerowych minimalnym kosztem -- dla pojedynczego interaktywnego formularza
    było to rozwiązanie adekwatne.

  \item \textbf{Świadome ograniczenie zakresu.} Rezygnacja z~pełnej rezerwacji
    online oraz z~tłumaczenia treści redakcyjnej pozwoliła dostarczyć kompletny,
    dopracowany zakres zamiast wielu niedokończonych funkcji.

  \item \textbf{Optymalizacja ma znaczenie już na małej skali.} Zastosowanie
    \texttt{prefetch\_related}/\texttt{select\_related} oraz cache dla danych rzadko
    zmienianych pokazało, jak prostymi środkami ogranicza się liczbę zapytań do bazy.

  \item \textbf{Kierunki rozbudowy.} Naturalne rozszerzenia to: wdrożenie
    produkcyjne za reverse proxy z~HTTPS, integracja z~zewnętrznym API (np. mapa
    okolicy lub pogoda na stronie atrakcji) oraz udostępnienie własnego API z~listą
    pokoi, pełne tłumaczenie treści (\texttt{django-modeltranslation}) i~podpięcie
    współdzielonego cache Redis (zależność jest już dołączona).
\end{enumerate}

\vspace{0.6cm}
\begin{center}
\begin{tabular}{lr}
\toprule
\textbf{Metryka} & \textbf{Wartość} \\
\midrule
Aplikacje Django       & 2 (core, content) \\
Strony publiczne       & 10 \\
Encje modelu danych    & 9 \\
Serwisy kontenerów     & 3 (web, db, mailhog) \\
Obsługiwane języki     & 2 (PL, DE) \\
\bottomrule
\end{tabular}
\end{center}

% ===================== DODATEK: SKOROWIDZ =====================
\newpage
\appendix
\section{Skorowidz zagadnień warstwy internetowej}
\label{sec:skorowidz}

Pomocniczy skorowidz wiążący zagadnienia omawiane na kursie z~miejscem ich realizacji
w~systemie oraz sekcją niniejszego dokumentu. Służy szybkiej nawigacji -- nie jest
listą kontrolną.

\begin{longtable}{p{0.30\linewidth}p{0.40\linewidth}c}
\toprule
\textbf{Zagadnienie} & \textbf{Realizacja w~systemie} & \textbf{Sekcja} \\
\midrule
\endhead
Framework MVC/MVT      & Django 4.2                            & \ref{sec:architektura} \\
Framework CSS          & Bootstrap 5.3 (CDN)                   & \ref{sec:frontend} \\
HTML / dziedziczenie   & \texttt{base.html} + partiale         & \ref{sec:frontend} \\
CSS własny             & \texttt{static/css/style.css}         & \ref{sec:frontend} \\
JavaScript             & \texttt{inquiry.js}, lightbox galerii & \ref{sec:frontend} \\
RWD                    & viewport + siatka + media queries     & \ref{sec:frontend} \\
Asynchroniczność       & HTMX (fragmenty HTML)                 & \ref{sec:frontend} \\
Routing / pretty URLs  & \texttt{urls.py}, \texttt{<slug>}     & \ref{sec:architektura} \\
ORM                    & modele + \texttt{prefetch/select\_related} & \ref{sec:orm} \\
Baza danych            & PostgreSQL 16                         & \ref{sec:dane} \\
Formularze             & \texttt{InquiryForm} + walidacja      & \ref{sec:backend} \\
Cache                  & \texttt{LocMemCache}, \texttt{@cache\_page} & \ref{sec:backend} \\
Poczta                 & \texttt{send\_mail} $\to$ MailHog     & \ref{sec:backend} \\
Uwierzytelnianie       & \texttt{LoginView}, \texttt{@login\_required} & \ref{sec:backend} \\
Logger                 & \texttt{LOGGING} + \texttt{logger.*}  & \ref{sec:backend} \\
Lokalizacja            & przełącznik PL/DE, \texttt{gettext}   & \ref{sec:i18n} \\
Dependency manager     & pip + \texttt{requirements.txt}       & \ref{sec:bezpieczenstwo} \\
Deployment             & Docker Compose                        & \ref{sec:uruchomienie} \\
\bottomrule
\end{longtable}

\end{document}
